%%%% IACR Transactions TEMPLATE %%%%
% This file shows how to use the iacrtrans class to write a paper.
% Written by Gaetan Leurent gaetan.leurent@inria.fr (2020)
% Public Domain (CC0)


%%%% 1. DOCUMENTCLASS %%%%
\documentclass[journal=tosc,submission]{iacrtrans}
%%%% NOTES:
% - Change "journal=tosc" to "journal=tches" if needed
% - Change "submission" to "final" for final version
% - Add "spthm" for LNCS-like theorems


%%%% 2. PACKAGES %%%%
\usepackage{lipsum} % Example package -- can be removed


%%%% 3. AUTHOR, INSTITUTE %%%%
\author{Jane Doe\inst{1,2} \and John Doe\inst{1}}
\institute{
  Institute A, City, Country, \email{jane@institute}
  \and
  Institute B, City, Country, \email{john@institute}
}
%%%% NOTES:
% - We need a city name for indexation purpose, even if it is redundant
%   (eg: University of Atlantis, Atlantis, Atlantis)
% - \inst{} can be omitted if there is a single institute,
%   or exactly one institute per author


%%%% 4. TITLE %%%%
\title{Evaluating ML-KEM at (much) higher security levels}
%%%% NOTES:
% - If the title is too long, or includes special macro, please
%   provide a "running title" as optional argument: \title[Short]{Long}
% - You can provide an optional subtitle with \subtitle.

\begin{document}

\maketitle


%%%% 5. KEYWORDS %%%%
\keywords{Something \and Something else}


%%%% 6. ABSTRACT %%%%
\begin{abstract}
\end{abstract}


%%%% 7. PAPER CONTENT %%%%
\section{Introduction}

After a multi-year, three-round public competition, NIST published its first post-quantum standards FIPS-203, FIPS-204 and FIPS-205 in August 2024,
specifying ML-KEM, ML-DSA and SLH-DSA respectively. Interestingly, four of the schemes in the final NIST PQC round --- Kyber (standardized as ML-KEM), Saber, NTRU LPRime, and FrodoKEM --- are closely related lattice KEMs that, setting aside their differing error mechanisms, span a spectrum of algebraic structure: from FrodoKEM's unstructured LWE, through Kyber's and Saber's cyclotomic modules, to NTRU~Prime's prime-degree fields.
This spectrum reflects different design principles, usually framed as structure vs.\ conservatism. ML-KEM's module structure gives small keys and fast NTT arithmetic, but is also seen as a latent attack surface, so most conservative designs shed it entirely. FrodoKEM, for example sits at the unstructured extreme, resting on plain LWE which is the reason bodies such as BSI and ANSSI point to FrodoKEM for long-term confidentiality.

It is useful to separate two kinds of conservatism that this framing bundles together, because they hedge against two independent (and, today, entirely hypothetical) classes of attack.
The first is the choice of assumption. FrodoKEM rests on unstructured LWE to guard against a \emph{structural break}, an attack that would exploit the ring's algebraic structure and would not apply to plain LWE.
The second is the security margin a scheme carries above its target level, which guards against a different threat, advances in the \emph{best known} attacks --- improved lattice reduction, dual attacks, sharper cost models --- that would erode the concrete security of every lattice scheme at its current parameters, FrodoKEM included.
The two hedges are distinct and non-substitutable: structural conservatism does nothing against the second class, and no margin protects against the first. Which of the two threats is the more likely to materialize, nobody can say.

Faced with that uncertainty, raising the margin is the historically standard response.
A concrete security estimate is not a proof but the cost of the best attack known at the time, and such estimates drift as cryptanalysis matures: the recommended size of an RSA modulus climbed from $512$ to $1024$ to $2048$ bits and beyond as factoring algorithms and available computation advanced, and estimates for lattice problems have likewise been revised repeatedly~\cite{TODO-lattice-estimates}.
A large margin is the standard insurance against such revision.

That insurance is costly without structure but cheap with it.
At matched security, the keys and ciphertexts of an unstructured scheme are already roughly an order of magnitude larger than a module scheme's~\cite{TODO-frodo-ring}; and because raising the module rank leaves the per-polynomial NTT arithmetic untouched, a module scheme can climb to a high margin that an unstructured scheme reaches only at far greater size.
Our work explores this direction, asking a question that has remained largely unexplored in the literature: how much security can a module KEM deliver at a fixed, modest bandwidth?

We answer this concretely with \textsf{Kaiburr}, a family of module KEMs obtained by scaling ML-KEM's module rank and revisiting its noise distribution.
Two parameter sets fit within the bandwidth budget of FrodoKEM-640-SHAKE (a $9.6$\,KB public key): \textsf{kaiburr6} ($k=18$) at a $6.9$\,KB public key and \textsf{kaiburr8} ($k=24$) at $9.2$\,KB.
Within that budget they reach $1162$- and $1557$-bit quantum security, roughly ten to thirteen times FrodoKEM-640's $119$ bits, at comparable decapsulation-failure probability ($<2^{-134}$) and $3$--$5\times$ fewer cycles.
The module structure thus converts a fixed bandwidth budget into a security margin no unstructured scheme can approach at comparable size.

\section{Main Result}
\label{sec:main}

%%%% 8. BILBIOGRAPHY %%%%
\bibliographystyle{alpha}
\bibliography{abbrev3,crypto,biblio}
%%%% NOTES
% - Download abbrev3.bib and crypto.bib from https://cryptobib.di.ens.fr/
% - Use bilbio.bib for additional references not in the cryptobib database.
%   If possible, take them from DBLP.

\end{document}
